四组点云都保留了建筑主体，但周围的杂点数量有明显差别。阈值较低时，建筑周围的零散点较多；提高阈值后，杂点减少，建筑外形更清楚。不过，参考视角的保留区域也随之缩小。0.95 的结果较干净，但部分表面变得稀疏，说明提高阈值也会损失一些细节。

表 3 与这一现象一致：阈值从 0.5 提高到 0.95 后，保留比例由 67.27\% 降至 56.22\%，融合点数减少约 22.0\%。其中，0.5 到 0.65 的变化较小，0.8 到 0.95 的保留比例下降更明显。重投影偏差也略有下降，表示留下的点在不同视角间更一致。由此可见，较高阈值能减少杂点，但不能同时保证点云更完整。
